\section{Conclusion}
\label{sec:6}

Cet article est parti d'une question simple pour un gestionnaire : face à une perturbation, quelle réaction mène au niveau de résilience visé, et après coup, la perte tenait-elle au choc, au réseau ou à la réaction ? Y répondre supposait des crises observées sous plusieurs gestions, ce qu'aucune donnée réelle ne fournit. La démarche proposée les produit par simulation, en rejouant chaque crise à l'identique sous plusieurs politiques, puis les mesure par une courbe de résilience à paramètres interprétables et en apprend un réseau bayésien dont l'effet de la politique est identifiable.

Sur 2 000 crises simulées sur le réseau ISOMORPH, la réaction agit sur la durée d'une crise plutôt que sur sa profondeur, et c'est le moment d'engagement des leviers, plus que leur nombre, qui fait leur rentabilité. Le réseau appris préconise une politique plus sobre dans environ une crise sur sept lorsque la durée du choc est connue, sans perte notable de résilience : l'information sur la durée d'un incident a donc une valeur décisionnelle propre.

Trois prolongements en découlent. Le premier est l'effet de la topologie : à perturbation et décision identiques, deux réseaux réagissent différemment, et le générateur permet d'étudier systématiquement cet effet sur plusieurs topologies, jusqu'à la reconfiguration du réseau elle-même comme levier de résilience. Le deuxième est économique : valoriser le coût des leviers et leur levée transformerait la préconisation en arbitrage explicite entre coût et résilience. Le troisième est empirique : confronter les préconisations à des crises réelles, même partiellement documentées.
